% GENERATED FILE. Edit canonical lessons, then run npm run generate:books.
% canonical-source-hash: fnv1a64:0c93b53f07206f68
% canonical-lessons: TE-C423-pratiphalam, TE-C423-vrtha, TE-C423-anavasaram, TE-C423-saripovu, TE-C423-viniyogincu

\chapter{Rewards and Waste}
\label{ch:te-rewards-and-waste}
\hlchaptermodality{423}

\begin{chapteropening}
\textbf{By the end of this chapter:} \emph{I can talk about rewards, waste and what is enough.}

Complete the last lesson of chapter 423: I can talk about rewards, waste and what is enough.
\end{chapteropening}

\section[\texorpdfstring{\te{ప్రతిఫలం} (pratiphalaṁ)}{ప్రతిఫలం (pratiphalaṁ)}]{\te{ప్రతిఫలం} (pratiphalaṁ) — a reward; a return}

\label{lesson:TE-C423-pratiphalam}



\begin{quote}
Before the new one: say the Telugu for a resolve, then the Telugu for to decide firmly.
\end{quote}

\subsection*{You'll want to know: \te{ప్రతిఫలం}}
\textbf{\te{ప్రతిఫలం}} — \emph{pratiphalaṁ} — \textquotedblleft{}a reward; a return\textquotedblright{}.

\begin{quote}
\te{కష్టానికి} \te{ప్రతిఫలం} \te{తప్పకుండా} \te{వస్తుంది}.
\end{quote}

\begin{quote}
\emph{kaṣṭāniki pratiphalaṁ tappakuṇḍā vastundi.}
\end{quote}

\begin{quote}
Hard work is sure to bring its reward.
\end{quote}

\te{ప్రతి}, in return, + \te{ఫలం}, fruit.

\begin{cousinweb}
\textbf{\te{ప్రతిఫలం}} (\emph{pratiphalaṁ}) is built from Sanskrit \emph{prati-}, \textquotedblleft{}back, in return\textquotedblright{}, and \emph{phala}, \textquotedblleft{}fruit, result\textquotedblright{} (Monier-Williams).
\end{cousinweb}

\begin{grammarlens}[title={What you've built}]
One more word for saying what happens if.
\end{grammarlens}

\subsection*{Your turn}
\begin{itemize}
\raggedright
  \item \emph{Say it:} \emph{pratiphalaṁ}
  \item \emph{Say it:} \emph{pratiphalaṁ}, once more
  \item \emph{Say it:} \emph{kaṣṭāniki pratiphalaṁ tappakuṇḍā vastundi.} — Hard work is sure to bring its reward.
  \item \emph{Recall:} say the Telugu for a resolve, then the Telugu for to decide firmly, then say \emph{pratiphalaṁ} again
\end{itemize}

\begin{tcolorbox}[breakable,colback=teal!4,colframe=teal!35!black,title={Before you move on}]
What does \te{ప్రతిఫలం} mean? (\textquotedblleft{}A reward; a return\textquotedblright{}.) Say it once more.
\end{tcolorbox}
\section[\texorpdfstring{\te{వృథా} (vṛthā)}{వృథా (vṛthā)}]{\te{వృథా} (vṛthā) — in vain; a waste}

\label{lesson:TE-C423-vrtha}



\begin{quote}
Before the new one: say the Telugu for to decide firmly, then the Telugu for a reward.
\end{quote}

\subsection*{You'll want to know: \te{వృథా}}
\textbf{\te{వృథా}} — \emph{vṛthā} — \textquotedblleft{}in vain; a waste\textquotedblright{}.

\begin{quote}
\te{అతని} \te{కష్టం} \te{అంతా} \te{వృథా} \te{అయింది}.
\end{quote}

\begin{quote}
\emph{atani kaṣṭaṁ antā vṛthā ayindi.}
\end{quote}

\begin{quote}
All his effort went to waste.
\end{quote}

\te{వృథా} \te{చేయు}, to waste: \te{సమయం} \te{వృథా} \te{చేయు}, to waste time.

\begin{cousinweb}
\textbf{\te{వృథా}} (\emph{vṛthā}) is borrowed from Sanskrit \emph{vṛthā}, \textquotedblleft{}at random; in vain, uselessly\textquotedblright{} (Monier-Williams).
\end{cousinweb}

\begin{grammarlens}[title={What you've built}]
One more word for saying what happens if.
\end{grammarlens}

\subsection*{Your turn}
\begin{itemize}
\raggedright
  \item \emph{Say it:} \emph{vṛthā}
  \item \emph{Say it:} \emph{vṛthā}, once more
  \item \emph{Say it:} \emph{atani kaṣṭaṁ antā vṛthā ayindi.} — All his effort went to waste.
  \item \emph{Recall:} say the Telugu for to decide firmly, then the Telugu for a reward, then say \emph{vṛthā} again
\end{itemize}

\begin{tcolorbox}[breakable,colback=teal!4,colframe=teal!35!black,title={Before you move on}]
What does \te{వృథా} mean? (\textquotedblleft{}In vain; a waste\textquotedblright{}.) Say it once more.
\end{tcolorbox}
\section[\texorpdfstring{\te{అనవసరం} (anavasaraṁ)}{అనవసరం (anavasaraṁ)}]{\te{అనవసరం} (anavasaraṁ) — needless; unnecessary}

\label{lesson:TE-C423-anavasaram}



\begin{quote}
Before the new one: say the Telugu for a reward, then the Telugu for in vain.
\end{quote}

\subsection*{You'll want to know: \te{అనవసరం}}
\textbf{\te{అనవసరం}} — \emph{anavasaraṁ} — \textquotedblleft{}needless; unnecessary\textquotedblright{}.

\begin{quote}
\te{ఆ} \te{ప్రయాణం} \te{అనవసరం}.
\end{quote}

\begin{quote}
\emph{ā prayāṇaṁ anavasaraṁ.}
\end{quote}

\begin{quote}
That trip is unnecessary.
\end{quote}

\te{అ}-, not, + \te{అవసరం}, need, which you know. \te{అనవసరంగా}, needlessly.

\begin{cousinweb}
\textbf{\te{అనవసరం}} (\emph{anavasaraṁ}) is borrowed from Sanskrit \emph{anavasara}, \textquotedblleft{}having no leisure; inopportune\textquotedblright{} (Monier-Williams), the negative of \emph{avasara}, \textquotedblleft{}occasion\textquotedblright{}, the word in \te{అవసరం}.
\end{cousinweb}

\begin{grammarlens}[title={What you've built}]
One more word for saying what happens if.
\end{grammarlens}

\subsection*{Your turn}
\begin{itemize}
\raggedright
  \item \emph{Say it:} \emph{anavasaraṁ}
  \item \emph{Say it:} \emph{anavasaraṁ}, once more
  \item \emph{Say it:} \emph{ā prayāṇaṁ anavasaraṁ.} — That trip is unnecessary.
  \item \emph{Recall:} say the Telugu for a reward, then the Telugu for in vain, then say \emph{anavasaraṁ} again
\end{itemize}

\begin{tcolorbox}[breakable,colback=teal!4,colframe=teal!35!black,title={Before you move on}]
What does \te{అనవసరం} mean? (\textquotedblleft{}Needless; unnecessary\textquotedblright{}.) Say it once more.
\end{tcolorbox}
\section[\texorpdfstring{\te{సరిపోవు} (saripōvu)}{సరిపోవు (saripōvu)}]{\te{సరిపోవు} (saripōvu) — to be enough; to suffice}

\label{lesson:TE-C423-saripovu}



\begin{quote}
Before the new one: say the Telugu for in vain, then the Telugu for needless.
\end{quote}

\subsection*{You'll want to know: \te{సరిపోవు}}
\textbf{\te{సరిపోవు}} — \emph{saripōvu} — \textquotedblleft{}to be enough; to suffice\textquotedblright{}.

\begin{quote}
\te{నా} \te{జీతం} \te{ఇంటి} \te{ఖర్చులకు} \te{సరిపోదు}.
\end{quote}

\begin{quote}
\emph{nā jītaṁ iṇṭi kharculaku saripōdu.}
\end{quote}

\begin{quote}
My salary is not enough for the household expenses.
\end{quote}

Past (it …): \te{సరిపోయింది}, \emph{saripōyindi}.

Its past is shown as \te{సరిపోయింది}, it was enough. Also to fit: \te{ఆ} \te{చొక్కా} \te{నాకు} \te{సరిపోయింది}, that shirt fitted me.

\begin{cousinweb}
\textbf{\te{సరిపోవు}} (\emph{saripōvu}) is built from \emph{sari}, \textquotedblleft{}right, equal\textquotedblright{}, and the native verb \emph{pōvu}, \textquotedblleft{}to go\textquotedblright{} (DEDR 4572).
\end{cousinweb}

\begin{grammarlens}[title={What you've built}]
One more word for saying what happens if.
\end{grammarlens}

\subsection*{Your turn}
\begin{itemize}
\raggedright
  \item \emph{Say it:} \emph{saripōvu}
  \item \emph{Say it:} \emph{saripōvu}, once more
  \item \emph{Say it:} \emph{nā jītaṁ iṇṭi kharculaku saripōdu.} — My salary is not enough for the household expenses.
  \item \emph{Recall:} say the Telugu for in vain, then the Telugu for needless, then say \emph{saripōvu} again
\end{itemize}

\begin{tcolorbox}[breakable,colback=teal!4,colframe=teal!35!black,title={Before you move on}]
What does \te{సరిపోవు} mean? (\textquotedblleft{}To be enough; to suffice\textquotedblright{}.) Say it once more.
\end{tcolorbox}
\section[\texorpdfstring{\te{వినియోగించు} (viniyōgiñcu)}{వినియోగించు (viniyōgiñcu)}]{\te{వినియోగించు} (viniyōgiñcu) — to make use of; to utilise}

\label{lesson:TE-C423-viniyogincu}



\begin{quote}
Before the new one: say the Telugu for needless, then the Telugu for to be enough.
\end{quote}

\subsection*{You'll want to know: \te{వినియోగించు}}
\textbf{\te{వినియోగించు}} — \emph{viniyōgiñcu} — \textquotedblleft{}to make use of; to utilise\textquotedblright{}.

\begin{quote}
\te{మనం} \te{నీటిని} \te{జాగ్రత్తగా} \te{వినియోగించాలి}.
\end{quote}

\begin{quote}
\emph{manaṁ nīṭini jāgrattagā viniyōgiñcāli.}
\end{quote}

\begin{quote}
We must use water carefully.
\end{quote}

Past (I …): \te{వినియోగించాను}, \emph{viniyōgiñcānu}.

Formal beside \te{ఉపయోగించు}, which you know.

\begin{cousinweb}
\textbf{\te{వినియోగించు}} (\emph{viniyōgiñcu}) is built on Sanskrit \emph{viniyoga}, \textquotedblleft{}distribution; employment, use, application\textquotedblright{} (Monier-Williams), with the Telugu verb ending -\te{ఇంచు}.
\end{cousinweb}

\begin{grammarlens}[title={What you've built}]
One more word for saying what happens if.
\end{grammarlens}

\subsection*{Your turn}
\begin{itemize}
\raggedright
  \item \emph{Say it:} \emph{viniyōgiñcu}
  \item \emph{Say it:} \emph{viniyōgiñcu}, once more
  \item \emph{Say it:} \emph{manaṁ nīṭini jāgrattagā viniyōgiñcāli.} — We must use water carefully.
  \item \emph{Recall:} say the Telugu for needless, then the Telugu for to be enough, then say \emph{viniyōgiñcu} again
\end{itemize}

\begin{tcolorbox}[breakable,colback=teal!4,colframe=teal!35!black,title={Before you move on}]
What does \te{వినియోగించు} mean? (\textquotedblleft{}To make use of; to utilise\textquotedblright{}.) Say it once more.
\end{tcolorbox}
